\documentclass[../../../include/open-logic-section]{subfiles}

\begin{document}

\olfileid{sth}{spine}{foundation}

\olsection{التأسيس}

لم يتم لنا المقصود \emph{كلّه}، وإنما قاربناه؛ وذلك أن~$\ZFminus$
لا تضمن أن تقع \emph{كل} مجموعة في بعض~${\OLId{latin-looped}{V}}_{\OLId{greek-variable}{alpha}}$، أي أن تكون كل
مجموعة قد نشأت في مرحلة من المراحل.

وقد يقال إن في وسعنا، من الوجه الرياضي المباشر، ألا \emph{نبالي}
بالمجموعات الخارجة عن الهرم؛ فإن وُجدت أغفلناها. غير أن تصورنا للمجموعة
إنما بُني على أن كل مجموعة تتكون في مرحلة ما (انظر~\stageshier{} في
\olref[z][story]{sec}). فوجب أن ننفي إمكان خروج مجموعة عن الهرم، ولا
سبيل إلى ذلك إلا بإضافة مسلّمة تجعل لكل مجموعة موضعًا فيه.

ولما كانت مجموعات~${\OLId{latin-looped}{V}}_{\OLId{greek-variable}{alpha}}$ هي مراحلنا، كان أقرب ما يلوح أن نتخذ
القول الآتي مسلّمة:

\begin{defish}
\emph{الانتظام.} $\forall {\OLId{latin-looped}{A}} \exists {\OLId{greek-variable}{alpha}}\, {\OLId{latin-looped}{A}} \subseteq {\OLId{latin-looped}{V}}_{\OLId{greek-variable}{alpha}}$
\end{defish}

وهذا مسلك مستقيم لا غبار عليه. إلا أننا سنعدل عنه، لعلة يأتي بيانها في
القسم التالي، إلى المسلّمة الآتية:

\begin{axiom}[التأسيس]
$(\forall {\OLId{latin-looped}{A}} \neq \emptyset)(\exists {\OLId{latin-looped}{B}} \in {\OLId{latin-looped}{A}}){\OLId{latin-looped}{A}} \cap {\OLId{latin-looped}{B}} = \emptyset$.
\end{axiom}

ومع شيء من العمل يتبين، داخل~$\ZFminus$، أن التأسيس يوجب الانتظام.
ونبدأ بالتعريف الآتي:
\begin{defn}
لكل مجموعة~${\OLId{latin-looped}{A}}$، لتكن:
\begin{align*}
	\text{cl}_{\OLMathNumeral{0}}({\OLId{latin-looped}{A}}) &= {\OLId{latin-looped}{A}},\\
	\text{cl}_{{\OLId{latin-ordinary}{n}}+{\OLMathNumeral{1}}}({\OLId{latin-looped}{A}}) &= \bigcup \text{cl}_{\OLId{latin-ordinary}{n}}({\OLId{latin-looped}{A}}),\\
	\text{trcl}({\OLId{latin-looped}{A}}) &= \bigcup_{{\OLId{latin-ordinary}{n}} < \omega} \text{cl}_{{\OLId{latin-ordinary}{n}}}({\OLId{latin-looped}{A}}).
\end{align*}
ونسمي~$\text{trcl}({\OLId{latin-looped}{A}})$ \emph{الإغلاق المتعدي} لـ${\OLId{latin-looped}{A}}$.
\end{defn}
\noindent
ووجه هذه التسمية ظاهر من القضية الآتية:
\begin{prop}\ollabel{subsetoftrcl}
${\OLId{latin-looped}{A}} \subseteq \trcl{{\OLId{latin-looped}{A}}}$ و$\trcl{{\OLId{latin-looped}{A}}}$ مجموعة متعدية.
\end{prop}

\begin{proof}
أما~${\OLId{latin-looped}{A}} = \text{cl}_{\OLMathNumeral{0}}({\OLId{latin-looped}{A}}) \subseteq \text{trcl}({\OLId{latin-looped}{A}})$ فظاهر. وإذا كان
${\OLId{latin-ordinary}{x}} \in {\OLId{latin-ordinary}{b}} \in \trcl{{\OLId{latin-looped}{A}}}$، كان~${\OLId{latin-ordinary}{b}} \in \text{cl}_{\OLId{latin-ordinary}{n}}({\OLId{latin-looped}{A}})$ لبعض~${\OLId{latin-ordinary}{n}}$؛ ومن ثم
${\OLId{latin-ordinary}{x}} \in \text{cl}_{{\OLId{latin-ordinary}{n}}+{\OLMathNumeral{1}}}({\OLId{latin-looped}{A}}) \subseteq \text{trcl}({\OLId{latin-looped}{A}})$.
\end{proof}

\begin{lem}\ollabel{lem:TransitiveWellFounded}
إذا كانت~${\OLId{latin-looped}{A}}$ مجموعة متعدية، وجد~${\OLId{greek-variable}{alpha}}$ يحقق
${\OLId{latin-looped}{A}} \subseteq {\OLId{latin-looped}{V}}_{\OLId{greek-variable}{alpha}}$.
\end{lem}

\begin{proof}
نستحضر تعريف «$\supstrict({\OLId{latin-looped}{X}})$» من
\olref[ordinals][opps]{defsupstrict}، ثم نعرّف مجموعتين:
\begin{align*}
	{\OLId{latin-looped}{D}}  &= \Setabs{{\OLId{latin-ordinary}{x}} \in {\OLId{latin-looped}{A}}}{\forall \delta\ {\OLId{latin-ordinary}{x}} \nsubseteq {\OLId{latin-looped}{V}}_\delta}\\
	{\OLId{greek-variable}{alpha}} &= \supstrict\Setabs{\delta}{(\exists {\OLId{latin-ordinary}{x}} \in {\OLId{latin-looped}{A}})
	({\OLId{latin-ordinary}{x}} \subseteq {\OLId{latin-looped}{V}}_\delta \land (\forall \gamma \in \delta){\OLId{latin-ordinary}{x}} \nsubseteq {\OLId{latin-looped}{V}}_\gamma)}
\end{align*}
ولنفرض أولًا~${\OLId{latin-looped}{D}} = \emptyset$. إذا كان~${\OLId{latin-ordinary}{x}} \in {\OLId{latin-looped}{A}}$، وجد~$\delta$ بحيث
${\OLId{latin-ordinary}{x}} \subseteq {\OLId{latin-looped}{V}}_\delta$؛ وبما أن الأعداد الترتيبية مرتبة ترتيبًا جيدًا،
جاز اختياره بحيث
$(\forall \gamma \in \delta){\OLId{latin-ordinary}{x}} \nsubseteq {\OLId{latin-looped}{V}}_\gamma$. فيكون
$\delta \in {\OLId{greek-variable}{alpha}}$، ويلزم~${\OLId{latin-ordinary}{x}} \in {\OLId{latin-looped}{V}}_{\OLId{greek-variable}{alpha}}$ من
\olref[spine][Valphabasic]{Valphabasicprops}. وإذن
${\OLId{latin-looped}{A}} \subseteq {\OLId{latin-looped}{V}}_{{\OLId{greek-variable}{alpha}}}$، وهو المطلوب.

فلم يبق إلا إثبات~${\OLId{latin-looped}{D}} = \emptyset$. نفرض خلافه طلبًا للمحال. يقضي
التأسيس بوجود~${\OLId{latin-looped}{B}} \in {\OLId{latin-looped}{D}} \subseteq {\OLId{latin-looped}{A}}$ بحيث~${\OLId{latin-looped}{D}} \cap {\OLId{latin-looped}{B}} = \emptyset$.
فإذا كان~${\OLId{latin-ordinary}{x}} \in {\OLId{latin-looped}{B}}$ كان~${\OLId{latin-ordinary}{x}} \in {\OLId{latin-looped}{A}}$ لتعدي~${\OLId{latin-looped}{A}}$؛ ولأن~${\OLId{latin-ordinary}{x}} \notin {\OLId{latin-looped}{D}}$،
ينتج~$\exists \delta\ {\OLId{latin-ordinary}{x}} \subseteq {\OLId{latin-looped}{V}}_\delta$. ضع الآن
\[
	\beta = \supstrict\Setabs{\delta}{(\exists {\OLId{latin-ordinary}{x}} \in {\OLId{latin-looped}{B}})({\OLId{latin-ordinary}{x}} \subseteq {\OLId{latin-looped}{V}}_\delta \land (\forall \gamma < \delta){\OLId{latin-ordinary}{x}} \nsubseteq {\OLId{latin-looped}{V}}_\gamma)}.
\]
فنستدل كما تقدم على~${\OLId{latin-looped}{B}} \subseteq {\OLId{latin-looped}{V}}_\beta$، وهذا يناقض كون~${\OLId{latin-looped}{B}} \in {\OLId{latin-looped}{D}}$.
\end{proof}

\begin{thm}\ollabel{zfentailsregularity}
يصدق الانتظام.
\end{thm}

\begin{proof}
خذ~${\OLId{latin-looped}{A}}$ كيفما اتفق. من~\olref{subsetoftrcl} نعلم أن
${\OLId{latin-looped}{A}} \subseteq \trcl{{\OLId{latin-looped}{A}}}$ وأن الأخيرة متعدية؛ فتوجد~${\OLId{greek-variable}{alpha}}$ بحيث
${\OLId{latin-looped}{A}} \subseteq \trcl{{\OLId{latin-looped}{A}}} \subseteq {\OLId{latin-looped}{V}}_{\OLId{greek-variable}{alpha}}$، بمقتضى
\olref{lem:TransitiveWellFounded}.
\end{proof}

فقد ثبت أن~$\ZFminus$ تقضي بـ
$\emph{التأسيس}\Rightarrow\emph{الانتظام}$. وسيأتي في
\olref[spine][rank]{zfminusregularityfoundation} أن~$\ZFminus$ تقضي
كذلك بـ$\emph{الانتظام}\Rightarrow\emph{التأسيس}$. فالتأسيس
والانتظام، مع افتراض~$\ZFminus$، متكافئان. وعلى هذا، بالنظر إلى
$\ZFminus$، يجوز تسويغ التأسيس بتكافئه مع الانتظام؛ أما الانتظام
نفسه فيسوّغه تصور~\stageshier{} مباشرة.

\end{document}
